\section{Complete source acquisition for the quadratic comparison}\label{sec:acquisition}

The preceding floor concerns the exact field. To state honestly what
a finite implementation of its sufficient comparison would require,
we retain the complete acquisition estimate from
\cite{OGMollifier,OGCorrection}. This is an analytical error bound;
no numerical certificate has been run.

\subsection{Arithmetic and radial source errors}

Let \(\Phi_{N_{\rm src}}\) be the first \(N_{\rm src}\) summands of
\eqref{eq:source} on \(u\ge0\), reflected evenly, and put
\[
 \widehat H(z)=\frac12\int_{-V}^{V}
          \Phi_{N_{\rm src}}(u)\ee^{\ii zu}\dd u.
\]
With \(B_0=5C_2(0)\), define
\[
 \mathcal Q=\int_\R\ee^{9|u|}\rho(u)^2\ee^{-\rho(u)}\dd u,\qquad
 \mathcal Q^{>V}=\int_{|u|>V}\ee^{9|u|}
                      \rho(u)^2\ee^{-\rho(u)}\dd u.
\]
Both integrals are finite by Section~\ref{sec:source}. Splitting the
difference into the arithmetic tail and the exact radial tail gives
\begin{equation}\label{eq:delta-H}
 \sup_{|\Im z|\le8}|H_0(z)-\widehat H(z)|
 \le\delta_H:=\frac12\{5C_2(N_{\rm src})\mathcal Q+
                          B_0\mathcal Q^{>V}\}.
\end{equation}
The terms tend to zero separately. Only the undifferentiated source
is used, so no smooth matching of the finite prefix at zero is assumed.

For \(|x-X|\le R\), \(0\le\Im z\le8\), the reciprocal-gamma
product \cite{DLMFReciprocalGammaOC} gives the compact multiplier bound
\begin{equation}\label{eq:multiplier}
 |D(z)A(s(z))|\le\mathcal B_A(X,R)
 :=16\pi^{9/4}C_A
 \exp\left\{\frac94+\frac{\pi^2[81+(|X|+R)^2]}{192}\right\}.
\end{equation}
To verify it, write \(w=s/2\) in
\[
 \Gamma(w)^{-1}=w\ee^{\gamma w}
        \prod_{n\ge1}(1+w/n)\ee^{-w/n}.
\]
The inequality
\(\log|1+w/n|\le\Re w/n+|w|^2/(2n^2)\) yields
\[
 |\Gamma(w)^{-1}|
 \le|w|\exp\{\gamma\Re w+\pi^2|w|^2/12\}.
\]
Now \(1/2\le\Re s\le9/2\), \(|s|\ge1/2\),
\(|s|^2\le[81+(|X|+R)^2]/4\), and \(\gamma<1\).
Insert these bounds into \eqref{eq:F} and use \(|A|\le C_A\).
With \(\widehat F_A=DA\widehat H\), the compact field error is at
most \(\mathcal B_A\delta_H\).
This bound is not used on the remote infinite boundary.

\subsection{Both exact boundary tails}

Let \(\mu_{X,R}\) be the restriction of \(\mu_X\) to \(|x-X|\le R\).
Equation~\eqref{eq:fractional} gives
\[
 |F_A(x)|\le C_A(2+|x|),\qquad
 |F_A(x+8\ii)|\le\zeta(9/2)^2.
\]
Writing \(B_X=2+|X|\), one obtains
\begin{align}\label{eq:quadratic-tail}
 \norm{F_A/c_A-1}_{L^2(\mu_X-\mu_{X,R})}^2
 &\le\tau_A(X,R),\nonumber\\
 \tau_A(X,R)
 &=\frac{\ee^{-aR}}{2a}\left\{
 \frac{C_A^2}{c_*^2}
 \left[(B_X+R)^2+\frac{2(B_X+R)}a+\frac2{a^2}\right]
 +\frac{\zeta(9/2)^4}{c_*^2}+2\right\}.
\end{align}
Indeed \(|v-1|^2\le2|v|^2+2\), and
\(K(u)\le\ee^{-a|u|}/8\).
Integrate the lower-boundary quadratic polynomial on both spatial
tails and then the upper-boundary constant on both tails.
The upper square requires the fourth power of \(\zeta(9/2)\).
All tails here belong to the exact field.

\subsection{Center acquisition and denominator propagation}

For an integer \(M_c\ge2\), the safe Euler logarithm has error
\[
 \left|\log\zeta(s_c)-\sum_{p^k\le M_c}\frac{p^{-ks_c}}k\right|
 \le\beta_{M_c}=\frac23M_c^{-3/2}.
\]
Keep the finite polynomial \(A(s_c)\) exactly and set
\[
 \widetilde c=q(s_c)A(s_c)
       \exp\left(\sum_{p^k\le M_c}\frac{p^{-ks_c}}k\right),\qquad
 |c_A-\widetilde c|\le
 \eta:=U_*(\ee^{\beta_{M_c}}-1).
\]
The modulus of the finite-log exponential is at most \(\zeta(5/2)\).
Future arithmetic evaluation error is added to \(\eta\).
No logarithm of \(A\) is used.

Let
\[
 \ell_c=|\widetilde c|-\eta>0,\qquad
 \mu_R=\mu_{X,R}(\partial{\rm strip})
       =\frac2\pi\arctan(\sinh(aR)),\qquad
 A_R=\norm{\widehat F_A/\widetilde c-1}_{L^2(\mu_{X,R})}.
\]
Then
\begin{equation}\label{eq:full-quadratic-acquisition}
 \epsilon_A(X)\le A_R+
 \frac{\eta}{\ell_c}(A_R+\sqrt{\mu_R})
 +\frac{\mathcal B_A\delta_H}{c_*}+\sqrt{\tau_A(X,R)}.
\end{equation}
For the proof, first split the complete \(L^2\) norm into compact
and remote pieces using the triangle inequality. On the compact
piece, add and subtract \(\widehat F_A/c_A\) and
\(\widehat F_A/\widetilde c\). The field error is at most
\(\mathcal B_A\delta_H/c_*\), since the restricted mass is at most one.
The denominator error is
\[
 \frac{|c_A-\widetilde c|}{|c_A|}
     \norm{\widehat F_A/\widetilde c}
 \le\frac{\eta}{\ell_c}(A_R+\sqrt{\mu_R}),
\]
which proves \eqref{eq:full-quadratic-acquisition}.

For each fixed \(X,A\) and a positive acquisition budget \(h\),
first choose \(R\) to make \(\sqrt{\tau_A}<h/3\).
For this \(R\), choose \(N_{\rm src},V\) so that
\(\mathcal B_A\delta_H/c_*<h/3\).
Finally choose the integer center cutoff and arithmetic precision so that
\(\eta<c_*/4\) and
\(2\eta\norm{\widehat F_A}_{L^2(\mu_{X,R})}/c_*^2<h/3\).
The exact denominator difference is then bounded by
\[
 \frac{\eta\norm{\widehat F_A}}{|c_A||\widetilde c|}
 \le\frac{4\eta\norm{\widehat F_A}}{3c_*^2}<h/3.
\]
This sharper expression, rather than requiring the conservative
surrogate in \eqref{eq:full-quadratic-acquisition} to meet the same
allocation, proves the ordered finite-acquisition assertion.
If the norm vanishes, the denominator error vanishes.
Every tuple can be treated this way, including a prescribed sequence
of growing heights or shrinking defect distances. This supplies no
common cutoff, affordable complexity bound, small residual value,
or evaluated exclusion.

